% ****************************************************************************************************
% Abstract / Front matter
% ****************************************************************************************************

\chapter*{Abstract}
\label{ch:abstract}
\addcontentsline{toc}{chapter}{Abstract}

This report describes a machine learning system that predicts the Emergency Severity Index (ESI) acuity level assigned to a patient at triage in an emergency department. The work began as an entry to the Triagegeist Kaggle competition sponsored by the Laitinen-Fredriksson Foundation. On examination of the provided dataset, I found that the chief complaint text had been synthetically generated with severity adjectives that encoded the label almost directly: a short keyword analysis and a verbatim train/test string overlap of $99.8\%$ together showed that a lookup table would score above $0.999$ on the competition metric. This discovery shifted the project from a competition entry to a small clinical-ML study on real emergency department records.

The real-data pipeline uses MIMIC-IV-ED, a credentialed dataset of $418{,}100$ emergency department visits from Beth Israel Deaconess Medical Center. The feature set combines $65$ tabular variables (raw vitals, engineered vitals including shock index and partial NEWS2 and qSOFA, medication-class flags, a leakage-safe Charlson comorbidity composite, and prior-visit history features) with $50$ principal components of fine-tuned Bio\_ClinicalBERT embeddings on the chief complaint text. The model is a blend of five gradient-boosted tree configurations combined by convex weight optimisation, with post-hoc per-class threshold multipliers for the final decision rule.

The out-of-fold macro F1 is $0.597$, Cohen's kappa is $0.701$, macro AUROC is $0.911$, and overall accuracy is $72.3\%$. The kappa sits close to the live-patient inter-rater reliability reported by Mirhaghi and colleagues ($\kappa = 0.694$) and in the substantial-agreement band of the Landis-Koch scale. The default argmax operating point produces a $15.2\%$ under-triage rate, well above the American College of Surgeons $5\%$ target. A sweep over training-time class-weighting and post-hoc threshold strategies produces a Pareto frontier of achievable (F1, under-triage) trade-offs; meeting the $5\%$ target requires accepting a macro F1 of roughly $0.44$--$0.46$. A substantive finding is that asymmetric training and post-hoc thresholds are substitute, not complementary, interventions: threshold optimisation after safety-weighted training largely undoes the training effect.

An intersectional fairness audit across age, sex, and race shows that age is the dominant disparity axis in the model's errors: patients aged $18$--$30$ are under-triaged at $1.71\times$ the rate of those aged $80+$, larger than the effects of race or sex in this cohort. The pattern is consistent with the clinical phenomenon of compensated shock, in which young adults maintain apparently normal vital signs despite developing critical illness. Age-adjusted vital-sign features close the disparity substantially at no meaningful F1 cost. I also document that the post-hoc safety threshold \emph{widens} the relative disparity between the best and worst subgroups even as it reduces every subgroup's absolute under-triage rate.

SHAP analysis of a representative base model reveals a strong embedding-dominance pattern: the first PCA component of the fine-tuned text embedding is individually more influential than any clinical feature, and sociodemographic variables (language, insurance) appear in the top ranking. The explainability story is therefore partial rather than complete. The report maps its work to the relevant articles of the EU AI Act (Regulation 2024/1689, Articles 9, 10, 13, and 14) and lists the specific gaps a deployed version of the system would still need to close.

Several limitations deserve explicit flagging. The work is single-institution, validated only on MIMIC-IV-ED. The race aggregation used in the fairness audit was incomplete and substantially undercounts minority populations; the race-level numbers should be treated as preliminary. Bootstrap confidence intervals, Brier/ECE/log-loss calibration metrics, SHAP interaction values, prospective outcome validation, and external-institution replication were all outside the scope of the project and are flagged as Tier 1 follow-ups. The contribution of this report is a retrospective methodological demonstration, not a deployable clinical system.
